\documentclass[reqno]{amsart}
\usepackage[T1]{fontenc}
\usepackage{amsmath,amssymb,amsthm}
\usepackage{xr-hyper}
\usepackage[hidelinks]{hyperref}
\externaldocument[gr-]{20261125T011000--galois-representations-black-box__workbook}
\externaldocument[bg-]{20261126T080700--black-box-galois-cohomology-of-number-fields__workbook}
\newcounter{chap}
% Chapter numbers, generated from 20261125T012000--grh-workbook__guide.org
% by scripts/make-chapters.py.  Do not edit: change the order of the
% index and rerun the script (./preview-site.sh and CI do this).
%
% Usage in a sheet, after \newcounter{chap}:
%   % Chapter numbers, generated from 20261125T012000--grh-workbook__guide.org
% by scripts/make-chapters.py.  Do not edit: change the order of the
% index and rerun the script (./preview-site.sh and CI do this).
%
% Usage in a sheet, after \newcounter{chap}:
%   % Chapter numbers, generated from 20261125T012000--grh-workbook__guide.org
% by scripts/make-chapters.py.  Do not edit: change the order of the
% index and rerun the script (./preview-site.sh and CI do this).
%
% Usage in a sheet, after \newcounter{chap}:
%   \input{chapters}\chapnumber{<sheet ID>}
\makeatletter
\providecommand\chapters@set[2]{\@namedef{chapters@#1}{#2}}
\chapters@set{20261125T001000}{1}% zeros-to-ramp-sums
\chapters@set{20261125T001100}{2}% inert-primes-and-estimate-b
\chapters@set{20261125T005000}{3}% galois-fourier-and-averaging
\chapters@set{20261125T005500}{4}% columns-from-a-sphere
\chapters@set{20261125T000000}{5}% slope-method
\chapters@set{20261126T030100}{6}% a-twist-for-each-prime
\chapters@set{20261125T003000}{7}% the-chain-and-the-square-root
\chapters@set{20261125T010000}{8}% rank-lemma-and-zero-estimate
\chapters@set{20261125T010500}{9}% poincare-series-and-near-orthogonality
\chapters@set{20261124T223000}{10}% elliptic-curves-for-the-proof
\chapters@set{20261125T009500}{11}% fock-model-and-theta-intertwiners
\chapters@set{20261124T223100}{12}% complex-multiplication-and-inert-reduction
\chapters@set{20261125T008000}{13}% one-cm-determinant
\chapters@set{20261124T214500}{14}% supersingular-theta-divisibility
\chapters@set{20261124T214600}{15}% the-isotropic-plane
\chapters@set{20261125T013700}{16}% pairings-on-a-formal-group-of-height-two
\chapters@set{20261125T013500}{17}% hecke-nondegeneracy-on-the-tree
\chapters@set{20261126T030200}{18}% the-bound-is-attained
\chapters@set{20261125T009000}{19}% kronecker-limit-and-chowla-selberg
\chapters@set{20261125T002000}{20}% definite-theta-lifts
\chapters@set{20261125T015500}{21}% siegel-weil-in-the-smallest-case
\chapters@set{20261125T002100}{22}% definite-theta-lifts-and-norms
\chapters@set{20261125T016500}{23}% hecke-operators-and-theta-lifts
\chapters@set{20261125T016700}{24}% divisibility-along-the-chain
\chapters@set{20261125T004000}{25}% denominators-and-sha
\chapters@set{20261125T004100}{26}% tate-twists-and-sha
\chapters@set{20261126T020430}{27}% the-minus-part-by-gauss-sums
\chapters@set{20261125T014100}{28}% congruence-modules
\chapters@set{20261125T014500}{29}% eisenstein-congruences-and-selmer-classes
\chapters@set{20261125T014700}{30}% pseudocharacters-and-the-length-bound
\chapters@set{20261125T011000}{31}% galois-representations-black-box
\chapters@set{20261126T030300}{32}% the-data
\chapters@set{20261126T020210}{33}% bessel-poincare-series-and-their-gram-operators
\chapters@set{20261126T030400}{34}% the-seesaw-identity
\chapters@set{20261126T030500}{35}% divisibility-of-the-composite-lift
\chapters@set{20261126T030600}{36}% theorem-a-short-forms
\chapters@set{20261126T030700}{37}% the-determinant-of-a-chain
\chapters@set{20261126T030800}{38}% the-cost-of-one-lift
\chapters@set{20261126T030900}{39}% denominators-of-projections-to-lifts
\chapters@set{20261126T031000}{40}% twists-with-few-congruences
\chapters@set{20261126T031100}{41}% theorem-c-and-the-main-inequality
\providecommand\chapnumber[1]{%
  \@ifundefined{chapters@#1}%
    {\PackageError{chapters}{Sheet #1 is not listed in the index}%
      {List it in 20261125T012000--grh-workbook__guide.org\MessageBreak
       and run scripts/make-chapters.py.}}%
    {\edef\thechap{\@nameuse{chapters@#1}}}}
\makeatother
\chapnumber{<sheet ID>}
\makeatletter
\providecommand\chapters@set[2]{\@namedef{chapters@#1}{#2}}
\chapters@set{20261125T001000}{1}% zeros-to-ramp-sums
\chapters@set{20261125T001100}{2}% inert-primes-and-estimate-b
\chapters@set{20261125T005000}{3}% galois-fourier-and-averaging
\chapters@set{20261125T005500}{4}% columns-from-a-sphere
\chapters@set{20261125T000000}{5}% slope-method
\chapters@set{20261126T030100}{6}% a-twist-for-each-prime
\chapters@set{20261125T003000}{7}% the-chain-and-the-square-root
\chapters@set{20261125T010000}{8}% rank-lemma-and-zero-estimate
\chapters@set{20261125T010500}{9}% poincare-series-and-near-orthogonality
\chapters@set{20261124T223000}{10}% elliptic-curves-for-the-proof
\chapters@set{20261125T009500}{11}% fock-model-and-theta-intertwiners
\chapters@set{20261124T223100}{12}% complex-multiplication-and-inert-reduction
\chapters@set{20261125T008000}{13}% one-cm-determinant
\chapters@set{20261124T214500}{14}% supersingular-theta-divisibility
\chapters@set{20261124T214600}{15}% the-isotropic-plane
\chapters@set{20261125T013700}{16}% pairings-on-a-formal-group-of-height-two
\chapters@set{20261125T013500}{17}% hecke-nondegeneracy-on-the-tree
\chapters@set{20261126T030200}{18}% the-bound-is-attained
\chapters@set{20261125T009000}{19}% kronecker-limit-and-chowla-selberg
\chapters@set{20261125T002000}{20}% definite-theta-lifts
\chapters@set{20261125T015500}{21}% siegel-weil-in-the-smallest-case
\chapters@set{20261125T002100}{22}% definite-theta-lifts-and-norms
\chapters@set{20261125T016500}{23}% hecke-operators-and-theta-lifts
\chapters@set{20261125T016700}{24}% divisibility-along-the-chain
\chapters@set{20261125T004000}{25}% denominators-and-sha
\chapters@set{20261125T004100}{26}% tate-twists-and-sha
\chapters@set{20261126T020430}{27}% the-minus-part-by-gauss-sums
\chapters@set{20261125T014100}{28}% congruence-modules
\chapters@set{20261125T014500}{29}% eisenstein-congruences-and-selmer-classes
\chapters@set{20261125T014700}{30}% pseudocharacters-and-the-length-bound
\chapters@set{20261125T011000}{31}% galois-representations-black-box
\chapters@set{20261126T030300}{32}% the-data
\chapters@set{20261126T020210}{33}% bessel-poincare-series-and-their-gram-operators
\chapters@set{20261126T030400}{34}% the-seesaw-identity
\chapters@set{20261126T030500}{35}% divisibility-of-the-composite-lift
\chapters@set{20261126T030600}{36}% theorem-a-short-forms
\chapters@set{20261126T030700}{37}% the-determinant-of-a-chain
\chapters@set{20261126T030800}{38}% the-cost-of-one-lift
\chapters@set{20261126T030900}{39}% denominators-of-projections-to-lifts
\chapters@set{20261126T031000}{40}% twists-with-few-congruences
\chapters@set{20261126T031100}{41}% theorem-c-and-the-main-inequality
\providecommand\chapnumber[1]{%
  \@ifundefined{chapters@#1}%
    {\PackageError{chapters}{Sheet #1 is not listed in the index}%
      {List it in 20261125T012000--grh-workbook__guide.org\MessageBreak
       and run scripts/make-chapters.py.}}%
    {\edef\thechap{\@nameuse{chapters@#1}}}}
\makeatother
\chapnumber{<sheet ID>}
\makeatletter
\providecommand\chapters@set[2]{\@namedef{chapters@#1}{#2}}
\chapters@set{20261125T001000}{1}% zeros-to-ramp-sums
\chapters@set{20261125T001100}{2}% inert-primes-and-estimate-b
\chapters@set{20261125T005000}{3}% galois-fourier-and-averaging
\chapters@set{20261125T005500}{4}% columns-from-a-sphere
\chapters@set{20261125T000000}{5}% slope-method
\chapters@set{20261126T030100}{6}% a-twist-for-each-prime
\chapters@set{20261125T003000}{7}% the-chain-and-the-square-root
\chapters@set{20261125T010000}{8}% rank-lemma-and-zero-estimate
\chapters@set{20261125T010500}{9}% poincare-series-and-near-orthogonality
\chapters@set{20261124T223000}{10}% elliptic-curves-for-the-proof
\chapters@set{20261125T009500}{11}% fock-model-and-theta-intertwiners
\chapters@set{20261124T223100}{12}% complex-multiplication-and-inert-reduction
\chapters@set{20261125T008000}{13}% one-cm-determinant
\chapters@set{20261124T214500}{14}% supersingular-theta-divisibility
\chapters@set{20261124T214600}{15}% the-isotropic-plane
\chapters@set{20261125T013700}{16}% pairings-on-a-formal-group-of-height-two
\chapters@set{20261125T013500}{17}% hecke-nondegeneracy-on-the-tree
\chapters@set{20261126T030200}{18}% the-bound-is-attained
\chapters@set{20261125T009000}{19}% kronecker-limit-and-chowla-selberg
\chapters@set{20261125T002000}{20}% definite-theta-lifts
\chapters@set{20261125T015500}{21}% siegel-weil-in-the-smallest-case
\chapters@set{20261125T002100}{22}% definite-theta-lifts-and-norms
\chapters@set{20261125T016500}{23}% hecke-operators-and-theta-lifts
\chapters@set{20261125T016700}{24}% divisibility-along-the-chain
\chapters@set{20261125T004000}{25}% denominators-and-sha
\chapters@set{20261125T004100}{26}% tate-twists-and-sha
\chapters@set{20261126T020430}{27}% the-minus-part-by-gauss-sums
\chapters@set{20261125T014100}{28}% congruence-modules
\chapters@set{20261125T014500}{29}% eisenstein-congruences-and-selmer-classes
\chapters@set{20261125T014700}{30}% pseudocharacters-and-the-length-bound
\chapters@set{20261125T011000}{31}% galois-representations-black-box
\chapters@set{20261126T030300}{32}% the-data
\chapters@set{20261126T020210}{33}% bessel-poincare-series-and-their-gram-operators
\chapters@set{20261126T030400}{34}% the-seesaw-identity
\chapters@set{20261126T030500}{35}% divisibility-of-the-composite-lift
\chapters@set{20261126T030600}{36}% theorem-a-short-forms
\chapters@set{20261126T030700}{37}% the-determinant-of-a-chain
\chapters@set{20261126T030800}{38}% the-cost-of-one-lift
\chapters@set{20261126T030900}{39}% denominators-of-projections-to-lifts
\chapters@set{20261126T031000}{40}% twists-with-few-congruences
\chapters@set{20261126T031100}{41}% theorem-c-and-the-main-inequality
\providecommand\chapnumber[1]{%
  \@ifundefined{chapters@#1}%
    {\PackageError{chapters}{Sheet #1 is not listed in the index}%
      {List it in 20261125T012000--grh-workbook__guide.org\MessageBreak
       and run scripts/make-chapters.py.}}%
    {\edef\thechap{\@nameuse{chapters@#1}}}}
\makeatother
\chapnumber{20261126T080800}
\renewcommand{\thesection}{\thechap.\arabic{section}}

\theoremstyle{definition}
\newtheorem{definition}{Definition}[section]
\newtheorem{problem}{Problem}[section]
\theoremstyle{plain}
\newtheorem{theorem}{Theorem}[section]
\theoremstyle{remark}
\newtheorem*{leanremark}{In the Lean formalization}

\newcommand{\Z}{\mathbb Z}
\newcommand{\Q}{\mathbb Q}
\newcommand{\C}{\mathbb C}
\newcommand{\Bmax}{B_{\max}}
\newcommand{\Bdr}{B_{\mathrm{dR}}}
\newcommand{\Bcris}{B_{\mathrm{cris}}}
\newcommand{\Ainf}{A_{\mathrm{inf}}}
\newcommand{\rk}{\operatorname{rk}}
\newcommand{\val}{\operatorname{val}}
\newcommand{\Eis}{\mathrm{Eis}}
\setlength{\emergencystretch}{2em}
\addtolength{\textheight}{3\baselineskip}

\title{Black box: crystalline lattices and the condition above \texorpdfstring{\(\ell\)}{ell}}
\author{}
\date{November 26, 2026}

\begin{document}
\maketitle
\refstepcounter{chap}\label{chap}

Chapter~\ref{gr-chap}, \emph{Black box: Galois representations of eigenforms},
introduces the representations attached to eigenforms. Here we consider
crystalline cocycles of lattices at places above \(\ell\). This chapter gives the rank of the cocycle module and a bound on the
index of its image after reduction. Tate twists illustrate both facts.

\paragraph{Prerequisites.}
Local fields, the cyclotomic character and continuous Galois cohomology,
including cup products and Kummer theory. Use the finite local duality
and Euler characteristic of Chapter~\ref{bg-chap},
Theorem~\ref{bg-thm:local-duality}. The module construction uses Witt vectors.

\section{Crystalline cocycles}\label{sec:cocycles}

We define the Hodge--Tate jumps and the crystalline cocycles of a lattice, then use the jumps to state the rank of the cocycle module.

Fix a prime \(\ell\) and a completion \(F_v\) of a number field \(F\) at \(v\mid\ell\), embedded
in \(\overline{\Q}_\ell\). Write \(G_{F_v}=\mathrm{Gal}(\overline{\Q}_\ell/F_v)\)
and \(\C_\ell=\widehat{\overline{\Q}_\ell}\).
A lattice \(L\) is finite free over \(\Z_\ell\) with continuous
\(G_{F_v}\)-action. Put \(V=L\otimes_{\Z_\ell}\Q_\ell\).
For \(e\in\Z\), write \(\Z_\ell(e)\) for the rank-one lattice with action through
\(\chi_{\mathrm{cyc}}^e\), where \(\chi_{\mathrm{cyc}}\) is the cyclotomic
character. Write \(L^*(1)=\operatorname{Hom}_{\Z_\ell}(L,\Z_\ell(1))\)
with the induced action. Tensors with period rings are over \(\Z_\ell\)
and carry the diagonal Galois action.

\begin{definition}\label{def:crystalline}
A lattice \(L\) is \emph{crystalline} if the fixed vectors of
\(G_{F_v}\) in \(\Bmax\otimes_{\Z_\ell}L\) span over \(\Bmax\).
It is \emph{potentially crystalline} if the restriction to the Galois
group of a finite extension of \(F_v\) is crystalline. Here \(\Bmax\)
is Colmez's crystalline period ring. Fontaine's ring \(\Bcris\) gives
the same conditions, using \(\varphi(\Bmax)\subseteq\Bcris\subseteq\Bmax\)
for the Galois-equivariant Frobenius \(\varphi\) \cite{Fo,Co}.
\end{definition}

\begin{definition}\label{def:jump}
A lattice is \emph{Hodge--Tate} if \(\C_\ell\otimes_{\Z_\ell} L\) is a
direct sum of twists \(\C_\ell(e)=\C_\ell\otimes_{\Z_\ell}\Z_\ell(e)\)
as a semilinear representation. A summand \(\C_\ell(e)\) contributes
one \emph{jump} \(u=-e\). In particular \(\Z_\ell(e)\) has jump \(-e\).
\end{definition}

Potentially crystalline lattices are Hodge--Tate and remain potentially
crystalline under Tate duality. The Tate twists are crystalline. A crystalline
lattice with all jumps zero is unramified, so inertia acts trivially
\cite{Fo}.

For a \(G_{F_v}\)-module write \(Z^1\) for the continuous cocycles and
\(B^1\) for the coboundaries \(g\mapsto gy-y\).
Write \(\mathrm{red}_n\) for pointwise reduction modulo \(\ell^n\).

\begin{definition}\label{def:Zf}
The module \(Z_f(L)\) consists of maps \(c\colon G_{F_v}\to L\) for which
there is a \(y\in\Bmax\otimes_{\Z_\ell}L\) with
\(1\otimes c(g)=gy-y\) for every \(g\in G_{F_v}\).
\end{definition}

\begin{theorem}[Rank of crystalline cocycles]\label{thm:rank}
Suppose \(L\) is potentially crystalline. Then \(Z_f(L)\) is a
saturated submodule of \(Z^1(F_v,L)\) and is finite free over \(\Z_\ell\).
If \(n^-\) is the number of negative jumps, counted with multiplicity, then
\[
  \rk_{\Z_\ell} Z_f(L)=\rk_{\Z_\ell} L+[F_v:\Q_\ell]n^- .
\]
\end{theorem}

Rank and finite generation are read from the proof \cite{GL}.
Write \(H^1_f(F_v,L)=Z_f(L)/B^1(F_v,L)\). The classical rational local
condition \(H^1_f(F_v,V)\) is the kernel of extension of scalars to
\(\Bcris\) on \(H^1(F_v,V)\). The map from \(H^1(F_v,L)\) to \(H^1(F_v,V)\)
pulls back the rational local condition to \(H^1_f(F_v,L)\) and induces
\(H^1_f(F_v,L)\otimes_{\Z_\ell}\Q_\ell\cong H^1_f(F_v,V)\).
Bloch and Kato give the rational condition and its inverse image
\cite[\S3]{BK}. The comparison of period rings in
Definition~\ref{def:crystalline} also gives the same kernel on cohomology
\cite{Co}. For the trivial representation the
rational condition consists of unramified homomorphisms, which factor
through \(G_{F_v}\) modulo inertia. For \(\Q_\ell(1)\) the rational
condition is the image under the Kummer map of the units of \(F_v\)
after \(\ell\)-adic completion and tensoring with \(\Q_\ell\)
\cite[Example~3.9]{BK}. Over \(\Q_\ell\), the completed unit group has
no \(\ell\)-torsion when \(\ell\) is odd.

The first problems compare the cocycles of the Tate twists with their coboundaries.

\begin{problem}\label{pr:tate-ranks}
For \(L=\Z_\ell(e)\) the ranks of \(Z_f(L)\), for \(e<0\), \(e=0\) and
\(e>0\), are \(1,1,1+[F_v:\Q_\ell]\), respectively. The corresponding
dimensions of \(H^1_f(F_v,\Q_\ell(e))\) are \(0,1,[F_v:\Q_\ell]\).
\end{problem}

\begin{problem}\label{pr:unramified}
The group \(Z_f(\Z_\ell)\) consists of the unramified homomorphisms,
with one free parameter given by the value on arithmetic Frobenius.
\end{problem}

\begin{problem}\label{pr:negative}
Take \(F_v=\Q_\ell\), \(\ell\) odd and \(e<0\). Put \(n_0=0\) if
\((\ell-1)\nmid e\) and \(n_0=1+\val_\ell(e)\) otherwise, with
\(\val_\ell(\ell)=1\). In a basis of \(\Z_\ell(e)\) the function
\[
  g\longmapsto \frac{\chi_{\mathrm{cyc}}(g)^e-1}{\ell^{n_0}}
\]
generates \(Z_f(\Z_\ell(e))\). The coboundaries have index \(\ell^{n_0}\).
For \((\ell,e)=(5,-4)\) and \((5,-20)\) the indices are \(5\) and \(25\).
\end{problem}

Reduction modulo \(\ell^n\) lets us count the two crystalline subgroups for a
lattice and its Tate dual.

\begin{problem}\label{pr:reduction}
For potentially crystalline \(L\) and \(n\ge 1\), the kernel of
\(\mathrm{red}_n\) on \(Z_f(L)\) is \(\ell^n Z_f(L)\).
Consequently \(\#\mathrm{red}_n Z_f(L)=\ell^{n\rk_{\Z_\ell} Z_f(L)}\).
\end{problem}

\begin{problem}\label{pr:dual-jump}
Let \(L\) be potentially crystalline and \(n\ge 1\). A jump \(u\) of
\(L\) corresponds to a jump \(-u-1\) of its Tate dual, with the same
multiplicity. Exactly one of the two integers is negative, and
\[
  \#\mathrm{red}_n Z_f(L)\,\#\mathrm{red}_n Z_f(L^*(1))
  =\lvert L/\ell^nL\rvert^{\,2+[F_v:\Q_\ell]} .
\]
For a Tate twist over \(\Q_\ell\) the product is \(\ell^{3n}\) for every \(e\).
\end{problem}

\section{Orthogonality}\label{sec:duality}

Local duality converts a restriction on cup products into a bound on an index
of crystalline cocycles.

\begin{theorem}[Crystalline local duality]\label{thm:duality}
Let \(L\) be potentially crystalline, let \(L'=L^*(1)\), and put
\(V'=L'\otimes_{\Z_\ell}\Q_\ell\).
In the local Tate pairing of \(H^1(F_v,L)\) with \(H^1(F_v,V'/L')\),
the subgroup \(H^1_f(F_v,L)\) and the image of \(H^1_f(F_v,V')\) are
exact annihilators.
\end{theorem}

The annihilator statement is the lattice form of the classical rational
duality of Bloch and Kato \cite[Proposition~3.8]{BK}.

For the index form keep \(L\) and \(L'\) as in Theorem~\ref{thm:duality}
and let \(n\ge 0\) be an integer. Let \(\mathcal Q\) be an additive
subgroup of maps \(G_{F_v}\to L'\), and \(\mathcal S\) an additive subgroup
of \(Z^1(F_v,L/\ell^n L)\). Suppose the cup product of every
\(s\in\mathcal S\cup\mathrm{red}_n Z_f(L)\) with \(\mathrm{red}_n c\) for
every \(c\in\mathcal Q\), evaluated in the group \(\mu_{\ell^n}\) of
\(\ell^n\)-th roots of unity, is the coboundary of a continuous \(1\)-cochain
in the discrete \(G_{F_v}\)-module \(\overline{\Q}_\ell^\times\). Then
\[
  [\mathcal S:\mathcal S\cap \mathrm{red}_n Z_f(L)]
  \ \mid\ [Z_f(L'):Z_f(L')\cap\mathcal Q] .
\]
An index of an infinite quotient is recorded as zero.
The index form is read from the proof \cite{GL}.

For the finite example take \(F_v=\Q_\ell\) and \(\ell\) odd. The two
modules are \(M=\Z/\ell\) with trivial action, and its dual \(\mu_\ell\).

\begin{problem}\label{pr:finite-counts}
Both first cohomology groups have order \(\ell^2\). The cocycle groups
of \(M\) and \(\mu_\ell\) have orders \(\ell^2\) and \(\ell^3\),
respectively. For \(\ell=5\) the cocycle orders are \(25\) and \(125\).
\end{problem}

\begin{problem}\label{pr:finite-annihilators}
The reductions modulo \(\ell\) of \(Z_f(\Z_\ell)\) and \(Z_f(\Z_\ell(1))\)
are exact annihilators for the pairing of \(Z^1(\Q_\ell,M)\) and
\(Z^1(\Q_\ell,\mu_\ell)\) induced by cup product and the local invariant.
Their orders are \(\ell\) and \(\ell^2\), respectively, and their images in
cohomology are the unramified classes and the Kummer classes of units.
For \(\ell=5\) the two orders are \(5\) and \(25\).
\end{problem}

\section{Modules over the Kisin ring}\label{sec:kisin}

The classical module has a Frobenius isomorphism after localization.
The chart frame gives an integral inclusion after multiplication by a
power of an Eisenstein element.

Write \(\C_\ell^\flat\) for the tilt, whose ring of integers is
\(\mathcal O_{\C_\ell^\flat}=\varprojlim_{y\mapsto y^\ell}
\mathcal O_{\C_\ell}/\ell\), and put \(\Ainf=W(\mathcal O_{\C_\ell^\flat})\).
Brackets denote Teichm\"uller lifts.

\begin{theorem}[Breuil--Kisin modules]\label{thm:kisin}
Choose a uniformizer \(\varpi_v\) of \(F_v\) with compatible \(\ell\)-power
roots. Put \(\mathfrak S=W(k)[[z]]\) for the residue field \(k\) of \(F_v\),
let \(\Eis\) be the Eisenstein polynomial of \(\varpi_v\), and let
\(\varphi\) be Witt-vector Frobenius extended by \(\varphi(z)=z^\ell\).
There is a fully faithful tensor functor from crystalline lattices \(L\)
to finite free \(\mathfrak S\)-modules \(\mathfrak M(L)\) with a
Frobenius isomorphism \(\varphi^*\mathfrak M(L)[1/\Eis]\cong
\mathfrak M(L)[1/\Eis]\). Extend scalars via the map sending \(z\) to the
Teichm\"uller lift of the chosen system of roots and then applying Witt
Frobenius to the resulting scalar map. Over \(W(\C_\ell^\flat)\)
the module is identified with \(L\otimes W(\C_\ell^\flat)\) compatibly
with Frobenius and the Galois group fixing the chosen roots.
\end{theorem}

The functor is classical \cite{Ki}\cite[Theorem~4.4]{BMS}. Its value on
\(\Z_\ell(1)\) is \(\mathfrak S\{1\}\), free of rank one with Frobenius
given in coordinates by a unit multiple of
\(y\mapsto \varphi(y)/\Eis\) \cite[Example~4.2 and Corollary~4.33]{BMS}.

For the chart form, choose a compatible sequence of primitive
\(\ell^j\)-th roots of unity and write
\(\epsilon=(1,\zeta_\ell,\zeta_{\ell^2},\ldots)\) in the tilt.
Put \(\tau=[\epsilon^{1/\ell}]-1\). Fontaine's map
\(\theta\colon\Ainf\to\mathcal O_{\C_\ell}\) sends \([\epsilon]\) to \(1\)
and \([\epsilon^{1/\ell}]\) to \(\zeta_\ell\). Its kernel is generated by
\(\sum_{i=0}^{\ell-1}[\epsilon^{1/\ell}]^i\) \cite[Example~3.16]{BMS}.
In the classical description, \(\Bdr^+\) is the completion of
\(\Ainf[1/\ell]\) at the kernel of the extension of \(\theta\), and
\(\Bdr\) is its fraction field \cite{Co}.

A Kummer chart has a ring \(\mathfrak S=W(k)[[z]]\subset\Ainf\), with
\(k\) a finite field of characteristic \(\ell\) and \(z\) a nonunit
Teichm\"uller lift. Its Eisenstein element \(\Eis\) is \(z^d\) times a
unit plus \(\ell\), with \(d\ge 1\), and generates \(\ker\theta\).
The subgroup \(\Gamma'\) of \(G_{\Q_\ell}\) fixing \(W(k)\) pointwise and
fixing \(\theta(z)\) lies in \(G_{F_v}\), has finite index prime to \(\ell\)
there, and contains a neighborhood of the identity in \(G_{F_v}\).
For every integer \(m\ge 0\) some \(g\in\Gamma'\) has
\(g(z)=[\epsilon]^m z\). Take \(L\) crystalline on \(\Gamma'\) and let
\(h\) be the least nonnegative bound for the absolute values of its
jumps on \(\Gamma'\). Write \(N_A\) for the submodule of \(\Ainf[1/\tau]\otimes L\)
consisting of vectors \(y\) with \(f(\varphi^j y)\in \Bdr^+\) for every
\(j\ge 1\) and every \(\Ainf\)-linear functional
\(f\colon W(\C_\ell^\flat)\otimes L\to\Bdr\) whose restriction to
\(1\otimes L\) is \(\Gamma'\)-equivariant. Frobenius acts on the Witt
vectors and fixes the lattice factor.

A consequence of the chart frame theorem \cite{GL} is the existence of a
free \(\mathfrak S\)-module \(N\) with a basis of \(\rk_{\Z_\ell}L\) vectors
that also freely generate \(N_A\) over \(\Ainf\).
In \(W(\C_\ell^\flat)\otimes L\) the inclusions are
\[
 \tau^h N_A\subseteq \Ainf\otimes L,\qquad
 \tau^h(\Ainf\otimes L)\subseteq N_A,\qquad
 \Eis^h\varphi N\subseteq N .
\]
The chart description and the frame consequence are taken as given.

\section{Limits with bounded jumps}\label{sec:limits}

The next two statements give conditions under which a limit is crystalline.

\begin{theorem}[Crystallinity from finite quotients]\label{thm:limit}
Let \(r\) be a positive integer. A lattice \(L\) is crystalline with jumps in \([-r,0]\)
if, for every \(n\ge 1\), \(L/\ell^n L\) is isomorphic as a \(G_{F_v}\)-module
to a quotient of two stable lattices in a crystalline representation
with jumps in \([-r,0]\).
\end{theorem}

The lattice theorem is Liu's \cite{Liu}, with the unramified base case due
to Berger \cite{Be}. The jumps use the sign of Definition~\ref{def:jump}.

For the trace form, fix a valuation ring \(A\) above \(\ell\) on the
algebraic closure \(\overline F\), with stabilizer \(D_A\) in
\(G_F=\operatorname{Gal}(\overline F/F)\). Fix
\(\iota_\ell\colon\C\simeq\overline{\Q}_\ell\).
A local datum is a field embedding \(j\colon\overline F\hookrightarrow
\overline{\Q}_\ell\) pulling back the valuation ring of
\(\overline{\Q}_\ell\) to \(A\), and a homomorphism
\(j_A\colon D_A\to G_{\Q_\ell}\) giving a topological isomorphism onto
the pointwise fixer of \(j(F)\), with \(j_A(s)(j(y))=j(sy)\) for \(s\in D_A\)
and \(y\in\overline F\).

Let \(r\) and \(\rho_{\eta i}\) be continuous finite-dimensional
representations of \(G_F\) over \(\overline{\Q}_\ell\). Here \(I\) is an
index set and \(i\) ranges over a finite set \(I_\eta\) for each
\(\eta\in I\). Fix
\(a_{\eta i}\in\overline{\Q}_\ell\). Suppose
\(\sum_{i\in I_\eta}a_{\eta i}\operatorname{tr}\rho_{\eta i}\)
converges uniformly on \(G_F\) to \(\operatorname{tr}r\) along a
nontrivial filter on \(I\), and \((y,y')\mapsto\operatorname{tr}(yy')\)
is nondegenerate on the \(\overline{\Q}_\ell\)-span of the matrices of \(r\).
For a local representation of \(D_A\) on a space \(V_{\mathrm{loc}}\) over
\(\overline{\Q}_\ell\), use Definition~\ref{def:crystalline} with \(V_{\mathrm{loc}}\)
in place of \(L\) and \(D_A\) acting on \(\Bmax\) through \(j_A\).
Assume each \(\rho_{\eta i}|_{D_A}\) is crystalline for every local datum.

For \(g\in G_{\Q_\ell}\), let \(\gamma\) be \(g\) followed by inclusion
into \(\C_\ell\). The multiplicity of \(u\in\Z\) at \(\gamma\) is the
dimension over \(\overline{\Q}_\ell\) of the space of \(\gamma\)-semilinear
maps \(f\colon V_{\mathrm{loc}}\to\C_\ell\) satisfying
\[
 f(sy)=\chi_{\mathrm{cyc}}(j_A(s))^{-u}j_A(s)(f(y))
\]
for each \(y\in V_{\mathrm{loc}}\) and all \(s\) sufficiently near the identity for that \(y\).
The multiplicity is recorded as zero if the dimension is infinite.

Assume either that \(I\) is finite or that there is one family of finite
multisets \(\mathcal J_\sigma\), indexed by embeddings \(\sigma\colon F\hookrightarrow\C\),
with the following property. For every \(\eta\in I\), \(i\in I_\eta\),
local datum with embedding \(j\), and \(g\in G_{\Q_\ell}\), the jump
multiplicities of \(\rho_{\eta i}|_{D_A}\) at \(\gamma\) are those of
\(\mathcal J_{\iota_\ell^{-1}\circ g^{-1}\circ j|_F}\).
Then \(r|_{D_A}\) is crystalline for every local datum. The trace
criterion is read from the proof \cite{GL}.

A sequence of Tate twists shows the need for control of jumps. For odd \(\ell\), the units of \(\Z_\ell\) are the product
\(\mu_{\ell-1}\times(1+\ell\Z_\ell)\), with \(\mu_{\ell-1}\) the
\((\ell-1)\)-st roots of unity. A continuous
character factoring through \(\chi_{\mathrm{cyc}}\) whose principal-unit
part is \(y\mapsto y^s=\exp(s\log y)\), for \(s\in\Z_\ell\), can be
Hodge--Tate only if \(s\in\Z\), by the rank-one Sen weight criterion \cite{Sen}.

\begin{problem}\label{pr:limit-exponent}
For \(n\ge 0\) put \(k_n=(\ell-1)(1+\ell^2+\cdots+\ell^{2n})\).
Then \(k_n\to -1/(\ell+1)\) in \(\Z_\ell\), and the jumps of
\(\Z_\ell(k_n)\) leave every bounded interval.
\end{problem}

\begin{problem}\label{pr:limit-character}
Take \(\ell\) odd and \(k_n\) as in Problem~\ref{pr:limit-exponent}.
The characters \(\chi_{\mathrm{cyc}}^{k_n}\) over \(\Q_\ell\) converge
uniformly to the character trivial on the torsion part of \(\Z_\ell^\times\)
and given by \(y\mapsto y^{-1/(\ell+1)}\) on the principal units.
The limit is not Hodge--Tate. At \(\ell=5\) its principal-unit exponent is
\(-1/6\).
\end{problem}

\section{A chart for a Tate twist}\label{sec:tate-chart}

The rank-one calculation compares the chart module with \(\Ainf\otimes L\).
In \(\Bdr^+\) the classical cyclotomic period \(t_{\mathrm{cyc}}=\log[\epsilon]\)
is a uniformizer with \(g(t_{\mathrm{cyc}})=\chi_{\mathrm{cyc}}(g)t_{\mathrm{cyc}}\), and
\((\Bdr)^{G_{F_v}}=F_v\) \cite{Fo,Co}.

\begin{problem}\label{pr:chart-period}
For \(j\ge 1\), \(\varphi^j\tau\) has order one in the discrete
valuation ring \(\Bdr^+\).
\end{problem}

\begin{problem}\label{pr:chart-tate}
On a chart over \(\Q_\ell\), take \(L=\Z_\ell(e)\) with \(e\ge 0\).
Then \(N_A=\tau^{-e}(\Ainf\otimes L)\).
\end{problem}

\begin{problem}\label{pr:chart-frob}
For the Tate twist of Problem~\ref{pr:chart-tate},
\(\Eis^e\varphi(N_A)\subseteq N_A\).
For \(e=1\) and a basis vector \(y\) of \(L\),
\(\Eis\varphi(\tau^{-1}y)\) is an \(\Ainf^\times\)-multiple of \(\tau^{-1}y\).
\end{problem}

\begin{leanremark}
For the crystalline conditions and the equation defining \(Z_f(L)\),
see \cite{GL}. The rank formula is
\path{crysCocycles.finrank_eq} in \path{BlochKato.lean}.
For the index see \path{relIndex_dvd_of_orthogonal} in
\path{CrystallineDuality.lean}.
\end{leanremark}

\paragraph{What was taken as given.}
The inputs are the classical crystalline period and jump facts, the
Bloch--Kato local conditions, the integral rank and index statements,
the Breuil--Kisin module and the chart frame, and the two limit statements.

{\small
\begin{thebibliography}{BMS}
\bibitem[Be]{Be} L.~Berger, \emph{Limites de repr\'esentations cristallines}, Compos.\ Math.\ \textbf{140} (2004), 1473--1498.
\bibitem[BMS]{BMS} B.~Bhatt, M.~Morrow and P.~Scholze, \emph{Integral \(p\)-adic Hodge theory}, Publ.\ Math.\ IH\'ES \textbf{128} (2018), 219--397.
\bibitem[BK]{BK} S.~Bloch and K.~Kato, \emph{\(L\)-functions and Tamagawa numbers of motives}, The Grothendieck Festschrift I, Progr.\ Math.\ 86, Birkh\"auser, 1990, 333--400.
\bibitem[Co]{Co} P.~Colmez, \emph{Th\'eorie d'Iwasawa des repr\'esentations de de Rham d'un corps local}, Ann.\ of Math.\ \textbf{148} (1998), 485--571.
\bibitem[Fo]{Fo} J.-M.~Fontaine, \emph{Repr\'esentations \(p\)-adiques semi-stables}, Ast\'erisque \textbf{223} (1994), 113--184.
\bibitem[Ki]{Ki} M.~Kisin, \emph{Crystalline representations and \(F\)-crystals}, Algebraic Geometry and Number Theory, Progr.\ Math.\ 253, Birkh\"auser, 2006, 459--496.
\bibitem[Liu]{Liu} T.~Liu, \emph{Torsion \(p\)-adic Galois representations and a conjecture of Fontaine}, Ann.\ Sci.\ \'Ec.\ Norm.\ Sup\'er.\ \textbf{40} (2007), 633--674.
\bibitem[Sen]{Sen} S.~Sen, \emph{Continuous cohomology and \(p\)-adic Galois representations}, Invent.\ Math.\ \textbf{62} (1980), 89--116.
\bibitem[GL]{GL} L.~Sernau, \emph{The Generalized Riemann Hypothesis for Dirichlet \(L\)-functions, in Lean~4}, \url{https://github.com/lsernau/grh-lean}, commit \texttt{72fead1}, readable edition \texttt{grh/}.
\end{thebibliography}
}

\section*{Hints}

\begin{description}
\item[\ref{pr:negative}] Use \(\Z_\ell^\times=\mu_{\ell-1}\times(1+\ell\Z_\ell)\) to find the least valuation of
\(\chi_{\mathrm{cyc}}(g)^e-1\).
\item[\ref{pr:chart-period}] Start with \(j=1\).
\item[\ref{pr:chart-tate}] Evaluate a chart functional on a basis of \(L\).
\end{description}

\end{document}
